\section{Conclusion}
\label{sec:conclusion}

We presented the first study of how self-critique reshapes internal representations in large language models. By extracting residual stream activations across 28 layers of \qwen at four pipeline stages on 200 \gsmk problems, we showed that self-critique induces significantly more representational drift than simple re-prompting ($p < 0.01$ at all layers after FDR correction). Critically, this drift is destructive: accuracy drops from 90.5\% to 84.0\%, linear probe separability decreases, and PCA analysis reveals disruption of fine-grained representational features. Self-critique is not shallow re-sampling---it is genuine restructuring that moves representations away from correct reasoning subspaces.

These findings have direct implications for the design of self-improvement agents. Systems that rely on intrinsic self-critique without external grounding risk degrading representations that were already well-calibrated. Our results suggest that simple re-prompting may be a safer alternative, and that representation-level interventions (\eg steering post-critique activations back toward pre-critique states) could potentially rescue the benefits of reflection while avoiding its costs.

Future work should extend this analysis to larger models, multiple domains, and causal intervention frameworks such as activation patching, to determine whether the destructive drift pattern is universal or model-specific, and whether it can be redirected toward constructive restructuring.
